% !TEX root = TCC - CD e IA.tex
% abstract em língua estrangeira (inglês)
%\setlength{\absparsep}{18pt} % ajusta o espaçamento dos parágrafos do resumo
\begin{otherlanguage*}{english}
This work proposes the development and validation of a hybrid text-mining methodology that serves as the engine of a decision-support tool for political communication teams. The approach integrates Natural Language Processing (NLP) techniques --- sentiment analysis, to identify polarities; topic modeling, to extract the central themes of the discussions --- and unsupervised learning (clustering), to group posts and profiles by similar characteristics. Unlike approaches that stop at describing public opinion, the results of these techniques are organized into an engagement funnel inspired by growth marketing and by the Consumers' Online Brand-Related Activities (COBRAs) framework, converting the analysis into actionable recommendations: what to produce, what to amplify, and what to abandon. The methodology distinguishes the discourse produced by the communication teams --- in captions and transcripts --- from the reaction expressed by the audience in the comments, and it is validated through a case study of Brazil's twenty-seven state governors. This research is expected to contribute a replicable analytical instrument for digital political marketing, overcoming the limitations of superficial engagement metrics by qualifying them through sentiment and organizing them in an actionable way.

\vspace{0.5cm}
\noindent
\textbf{Keywords}: Digital Political Marketing, Sentiment Analysis, Topic Modeling, Engagement (COBRAs), Growth, Data Science.
\end{otherlanguage*}